\chapter{Realización Palatini–Cartan–Holst del parámetro de Barbero–Immirzi y transducción física del espectro de área}
\label{chap:parte-iv-68}

La derivación estructural del parámetro y del espectro de área antecede a este capítulo. Aquí se construyen tétrada, conexión, curvatura y acción de primer orden, y sólo entonces se especializa el parámetro de Holst como realización física del lector areal.

\section{Curvatura, torsión y acción Einstein--Cartan--Holst}
\label{chap:gravedad-cartan-holst}

La estructura de transporte HMT distingue la variable de memoria, la
holonomía y el defecto de cierre. Su realización gravitatoria se formula
mediante una tétrada y una conexión. Esta sección establece la geometría de
primer orden y especifica el tipo del mapa entre ésta y el espacio de estados
TPK\@.

\subsection{Identidad discreta de Gauss–Stokes y transporte de incidencia a la frontera}

Sea \(K\) un complejo celular, sea \(W\) un grupo abeliano y sea
\(\Omega\in C^1(K;W)\) una cocadena. Denotamos por
\(\delta:C^1(K;W)\to C^2(K;W)\) el coborde celular y por
\(\langle\,\cdot\,,\,\cdot\,\rangle\) el emparejamiento entre cocadenas con
valores en \(W\) y cadenas celulares enteras.

\begin{theorem}[Identidad celular universal de Stokes]
Para toda \(2\)-cadena finita \(z\in C_2(K;\mathbb Z)\),
\[
 \boxed{\langle\delta\Omega,z\rangle
 =\langle\Omega,\partial z\rangle.}
\]
\end{theorem}
\begin{proof}
Es la identidad que define el coborde como operador adjunto de la frontera
celular. La finitud de \(z\) garantiza que ambos emparejamientos son sumas
finitas en \(W\).
\end{proof}

\begin{corollary}[Gauss--Stokes sobre una pseudovariedad celular]
Sea \(K\) una pseudovariedad celular regular, bidimensional, finita y compacta,
coherentemente orientada, posiblemente con borde. Si \(z_K\) es la suma de
sus \(2\)-celdas con la orientación coherente, entonces
\[
 \sum_{f\in K^{(2)}}(\delta\Omega)(f)
 =\sum_e[\partial z_K:e]\,\Omega(e).
\]
En particular, la regularidad permite identificar las incidencias no
canceladas de \(\partial z_K\) con las aristas del borde geométrico.
\end{corollary}
\begin{proof}
Se aplica el teorema a \(z_K\). La orientación coherente hace que cada arista
interior aparezca en \(\partial z_K\) dos veces con coeficientes opuestos;
las incidencias no canceladas forman la cadena de borde.
\end{proof}

La versión no abeliana no se obtiene sustituyendo sumas por productos sin más.
En esta obra se utiliza únicamente cuando se han fijado transportes a una base
común, un orden de las holonomías y la conjugación correspondiente. Fuera de
ese subdominio no se afirma una identidad de Stokes no abeliana.

\subsection{Capacidad areal de la celda qutrit: \(81\log 3\) y transporte colectivo de incidencia}

\begin{theorem}[Corte mínimo de un cubo discreto]
En el grafo cúbico \(N\times N\times N\) con capacidad unitaria, todo corte
que separa dos caras opuestas tiene capacidad al menos \(N^2\), y existe uno
de capacidad \(N^2\).
\end{theorem}
\begin{proof}
Hay \(N^2\) caminos disjuntos en aristas que atraviesan el cubo en la dirección
normal a las dos caras. Todo corte debe interceptarlos, luego su capacidad es
al menos \(N^2\). Una capa transversal contiene exactamente \(N^2\) aristas y
realiza la cota.
\end{proof}

Para \(N=9\), el corte vale 81 y
\[
 9^3=729=27^2.
\]
La última igualdad permite reagrupar seis trits como un volumen \(9^3\) o una
cuadrícula de cardinalidad \(27^2\). Es una biyección de conjuntos; no es un
isomorfismo entre
\((\ZZ/9\ZZ)^3\) y \((\ZZ/27\ZZ)^2\), cuyos exponentes son distintos.

\subsection{Tétrada, conexión y \(2\) defectos}

Sean \(e^I\) una tétrada y \(\omega^{IJ}=-\omega^{JI}\) una conexión de
Lorentz. Definimos
\[
 T_{\rm tor}^I=D_\omega e^I
 =\dd e^I+\omega^I{}_J\wedge e^J,
\]
\[
 F^{IJ}=\dd\omega^{IJ}+\omega^I{}_K\wedge\omega^{KJ},
 \qquad
 g=\eta_{IJ}e^I\otimes e^J.
\]
La descomposición
\[
 \omega^{IJ}=\mathring\omega^{IJ}(e)+K^{IJ}
\]
separa la conexión de Levi--Civita y la contorsión. Entonces
\[
 T_{\rm tor}^I=K^I{}_J\wedge e^J.
\]
Curvatura y torsión son defectos diferentes: \(F\) mide el cierre del
transporte de vectores; \(T_{\rm tor}\), el cierre de la tétrada. La memoria
\(T\) del principio de Ambivalencia tampoco es torsión geométrica: sólo un
funtor de realización puede relacionarlas.

\subsection{Acción de 1.er orden}

Sea \(M\) una variedad lisa orientada de dimensión cuatro. Si
\(\partial M\ne\varnothing\), suponemos variaciones de soporte compacto en el
interior o, equivalentemente para el problema variacional considerado, un
término de borde y condiciones de contorno que anulen la contribución
fronteriza. Cuando \(\Psi\) contenga campos espinoriales, suponemos además una
estructura de espín y campos compatibles con ella. Sean
\[
 \Sigma^{IJ}=e^I\wedge e^J,
 \qquad
 (\star X)^{IJ}=\frac12\epsilon^{IJ}{}_{KL}X^{KL},
 \qquad
 \kappa=\frac{8\pi G}{c^4},
 \qquad
 \gamma\in\mathbb R\setminus\{0\}.
\]
La acción Palatini--Cartan--Holst es
\[
\boxed{
 S[e,\omega,\Psi]
 =\frac1{2\kappa c}\int_M
 \Sigma_{IJ}\wedge\left(\star F^{IJ}+\frac1\gamma F^{IJ}\right)
 -\frac\Lambda{\kappa c}\int_M\operatorname{vol}_e
 +S_{\rm m}[e,\omega,\Psi].}
\]
La variación de esta acción determina las correcciones torsionales
\cite{Sciama1964,Kibble1961,Hehl1976,Holst1996}.

Definimos la corriente de espín por
\[
 \delta_\omega S_{\rm m}
 =-\frac12\int_M\delta\omega^{IJ}\wedge\sigma_{IJ}.
\]

\begin{theorem}[Ecuación de Cartan--Holst]
Bajo las hipótesis geométricas y de borde anteriores, la variación respecto
de \(\omega\) da
\[
 D_\omega\!\left[(\star+\gamma^{-1}I)\Sigma^{IJ}\right]
 =\kappa c\,\sigma^{IJ},
\]
con la normalización global fijada por la definición anterior de
\(\sigma^{IJ}\).
\end{theorem}
\begin{proof}
Se usa \(\delta F^{IJ}=D_\omega\delta\omega^{IJ}\), se integra por partes y
se iguala a cero el coeficiente de la variación arbitraria
\(\delta\omega^{IJ}\). Con la convención de signo fijada en la definición de
\(\sigma^{IJ}\), el término de materia pasa al miembro derecho con signo
positivo. El factor \(1/(2\kappa c)\) produce el coeficiente
\(\kappa c\). El término cosmológico no depende de \(\omega\), y las
condiciones de borde eliminan el término generado por la integración por
partes.
\end{proof}

En firma lorentziana \(\star^2=-I\) sobre bivectores y, para
\(\gamma\in\RR\setminus\{0\}\),
\[
 (\star+\gamma^{-1}I)^{-1}
 =\frac{\gamma^2}{1+\gamma^2}(-\star+\gamma^{-1}I).
\]

\begin{corollary}[Desacoplamiento torsional]
Para una co-tétrada no degenerada y
\(\gamma\in\mathbb R\setminus\{0\}\),
\[
 \sigma^{IJ}=0\Longrightarrow T_{\rm tor}^I=0
 \Longrightarrow\omega=\mathring\omega(e).
\]
\end{corollary}

Para hacer explícita la normalización, definimos las fuentes normalizadas por
la acción mediante
\[
 \delta_g S_{\rm m}
 =-\frac12\int_M\operatorname{vol}_e\,
 \mathcal T^{S}_{\mu\nu}\,\delta g^{\mu\nu}.
\]
Cuando las coordenadas de la tétrada tienen dimensión de longitud, el tensor
energía--momento físico convencional es
\(T^{\rm phys}_{\mu\nu}:=c\,\mathcal T^{S}_{\mu\nu}\). Denotamos de igual
modo por \(\mathcal U^{S,\rm spin}\) la contribución normalizada por la
acción y por \(U^{\rm phys,\rm spin}:=c\,\mathcal U^{S,\rm spin}\) su
convención física. Al eliminar algebraicamente la conexión, la ecuación
métrica adopta entonces las dos formas equivalentes
\[
 G_{\mu\nu}+\Lambda g_{\mu\nu}
 =\kappa c\,
 \bigl(\mathcal T_{\mu\nu}^{S,\rm LC}
       +\mathcal U_{\mu\nu}^{S,\rm spin}\bigr)
 =\kappa\,
 \bigl(T_{\mu\nu}^{\rm phys,LC}
       +U_{\mu\nu}^{\rm phys,spin}\bigr),
\]
donde la contribución de espín es cuadrática en la corriente correspondiente
y depende del acoplamiento de materia. El coeficiente de esa contribución no
se congela sin especificar la acción fermiónica y las convenciones de
dualidad.

\subsection{Identidad de Bianchi con torsión y balance de flujo en la conexión de Cartan}

Para toda conexión,
\[
 D_\omega T_{\rm tor}^I=F^I{}_J\wedge e^J,
 \qquad
 D_\omega F^{IJ}=0.
\]
Después de eliminar torsión,
\[
\nabla^\mu(T_{\mu\nu}^{\rm phys,LC}
          +U_{\mu\nu}^{\rm phys,spin})=0.
\]
Estas identidades son pruebas de consistencia obligatorias: cualquier término
gravitatorio HMT debe surgir de una acción compatible o satisfacer el balance
de Noether correspondiente.

\subsection{Mapa de realización y especialización del parámetro de Holst}

La inserción definida en el \cref{mdg:chap:barbero} toma
\[
 \gamma=\Gamma_{6\to8}.
\]
Una vez realizada esta sustitución, todas las consecuencias variacionales
anteriores son teoremas de la acción especializada. La selección física de
esta acción por el transporte HMT se expresa mediante un mapa de realización
\[
 \mathfrak R_{\rm grav}:\mathcal X_{\rm TPK}
 \dashrightarrow\{(e,\omega,\Psi)\}
\]
dotado de tres propiedades: entrelazar el transporte TPK con el de
\(\omega\), identificar las clases de holonomía y suministrar las escalas de
\(G\) y \(\Lambda\). Estas propiedades fijan el dominio preciso de una
realización gravitatoria del formalismo.

Para co-tétrada no degenerada, \(\gamma\) real no nulo y materia sin corriente
de espín, el sector torsional se reduce a la formulación de primer orden de
relatividad general con la misma constante cosmológica, sujeto a las
condiciones de borde ya declaradas. En un fluido de espín, el escalado
\(s^2\propto a^{-6}\) puede producir un rebote Einstein--Cartan; obtener su
densidad y sus condiciones iniciales a partir del estado TPK es una cuestión
distinta de la consistencia variacional de la acción.


\section{Retorno nonádico, grupoide de memoria y condición de entrelazamiento de Cartan}

La jerarquía \(27\to54\to108\), la cobertura acumulativa dirigida y el
grupoide de memoria suministran la condición exacta bajo la cual el retorno
del TPK admite una conexión de Cartan compatible. Este puente antecede a la
especialización Palatini--Cartan--Holst y conserva íntegros su dominio y sus
criterios de refutación.

\begingroup
\let\chapter\section
\let\section\subsection
\let\subsection\subsubsection
\let\subsubsection\paragraph
\let\clearpage\relax
\let\cleardoublepage\relax
\chapter{Retorno nonádico, memoria y holonomía de Cartan}
\label{chap:puente-retorno-holonomia-cartan}
\index{retorno!27--54--108}\index{holonomía!TPK--Cartan}
\index{memoria!distinción respecto de torsión}

La realización gravitatoria no comienza identificando una variable discreta
con una magnitud geométrica. Comienza por distinguir tres niveles: el retorno
del transporte observable, la memoria que diferencia historias con igual
proyección y la holonomía de una conexión. El puente que se formula aquí
recibe los retornos y cociclos ya demostrados en HMT; no los reconstruye
desde la acción de Holst, desde una masa ni desde una constante metrológica.

\section{Jerarquía exacta de los retornos}
\label{sec:jerarquia-retornos-27-54-108}

Sea \(F_{\rm obs}\) la cronología observable definida en el núcleo HMT y sea
\[
 \mathcal K_{27}:=F_{\rm obs}^{27}.
\]
El símbolo \(\mathcal K_{27}\) designa exclusivamente la composición de
veintisiete actualizaciones del lazo posicional y orientado; no designa una
ventana genérica, una traza de \(108\) pasos ni el estado TPK enriquecido
completo. Denotemos por \(p\) la proyección a las posiciones de los dos
cursores y por \(\mathsf J_{\rm or}\) la inversión simultánea de sus
orientaciones. La demostración primaria se encuentra en
\cref{prop:retorno-fobs-identidad}; sus consecuencias tipadas son
\[
 p(\mathcal K_{27}x)=p(x),\qquad
 o(\mathcal K_{27}x)=\mathsf J_{\rm or}\,o(x),
 \qquad \mathsf J_{\rm or}^{\,2}=I,
\]
\[
 F_{\rm obs}^{54}:
 \text{restaura posiciones y orientaciones},
 \qquad
 F_{\rm obs}^{108}=I:
 \text{restaura además la fase observable}.
\]
Por tanto, retorno de posición, retorno de orientación y retorno de fase son
afirmaciones diferentes. Ninguna de ellas implica por sí sola el retorno de
todos los campos de una extensión enriquecida.

El cociente de modos y los cociclos de orientación se establecen en
\cref{sec:cociclos-modo-orientacion,thm:descenso-necesario-fav}. En el orden
areal--volumétrico,
\[
 N_{27}=9F_{\rm av},
 \qquad
 F_{\rm av}=
 \begin{pmatrix}0&1\\1&1\end{pmatrix}.
\]
Si
\[
 \Omega_{\rm sw}(\gamma)=\sum_{a\subset\gamma}\omega_{\rm sw}(a),
 \qquad
 S(\gamma)=\sum_{a\subset\gamma}|\omega_{\rm sw}(a)|,
 \qquad
 G_{\rm or}(\gamma)=\sum_{a\subset\gamma}g(a),
\]
entonces un bloque de veintisiete pasos contiene tres vueltas nonádicas y una
inversión de orientación:
\[
 \Omega_{\rm sw}(\mathcal K_{27})=0,\qquad
 S(\mathcal K_{27})=18,\qquad
 G_{\rm or}(\mathcal K_{27})=1.
\]
En el retorno observable de \(108\) pasos,
\[
 \Omega_{\rm sw}(C_{108})=0,\qquad
 S(C_{108})=72,\qquad
 G_{\rm or}(C_{108})=4.
\]
Así, para el funcional aditivo declarado
\[
 \mathcal A_\ast(\gamma)
 :=G_{\rm or}(\gamma)+\lambda_\ast S(\gamma),
\]
se obtienen, sin introducir una escala dimensional,
\[
 \mathcal A_\ast(\mathcal K_{27})=1+18\lambda_\ast,
 \qquad
 \mathcal A_\ast(C_{108})=4+72\lambda_\ast.
\]
Estos valores son conteos del registro. Su posterior uso térmico o
gravitatorio no altera su procedencia.

\subsection{Sistema dirigido de coberturas, refinamiento y límite de la conexión de Cartan}
\label{sec:cobertura-acumulativa-dirigida}

Sea \(\mathcal O_{108}\) la órbita observable y, para cada arista cronológica
\(e\), sea
\[
 c(e)=\bigl(g(e),s(e)\bigr)\in\mathbb N^2
\]
el par formado por los indicadores de inversión de orientación y de cambio
efectivo de hoja. Para una historia dirigida
\(\gamma=e_1\cdots e_n\), se define
\[
 c(\gamma)=\sum_{j=1}^{n}c(e_j).
\]
La ley de concatenación toma la forma
\[
 c_{m+n}(x)
 =c_m(x)+c_n\!\left(F_{\rm obs}^{m}x\right).
\]
Los conteos anteriores expresan, de manera uniforme sobre la órbita,
\[
 c_{27}(x)=(1,18),
 \qquad
 c_{108}(x)=(4,72).
\]

\begin{theorem}[Cobertura acumulativa de la cronología observable]
\label{thm:cobertura-acumulativa-dirigida}
Sobre
\[
 \widetilde{\mathcal O}_{108}^{+}
 :=\mathcal O_{108}\times\mathbb N^2
\]
defínase
\[
 \widetilde F^{+}(x,m)
 :=\bigl(F_{\rm obs}x,m+c(x\to F_{\rm obs}x)\bigr).
\]
Entonces el transporte de veintisiete pasos satisface
\[
 \widetilde{\mathcal K}_{27}^{+}(x,m)
 =\bigl(\mathcal K_{27}x,m+(1,18)\bigr),
\]
y
\[
 \left(\widetilde{\mathcal K}_{27}^{+}\right)^4(x,m)
 =\bigl(x,m+(4,72)\bigr)\neq(x,m).
\]
Por tanto, su proyección observable tiene orden cuatro, mientras que la
acción acumulativa del monoide cronológico posee orden semigrupal infinito.
\end{theorem}

\begin{proof}
La ley cocíclica y la constancia de \(c_{27}\) dan
\[
 c_{108}(x)
 =\sum_{j=0}^{3}c_{27}(\mathcal K_{27}^{j}x)
 =4(1,18)=(4,72).
\]
La primera coordenada retorna por
\(\mathcal K_{27}^{4}=I\); la segunda no retorna. Tras \(q>0\) iteraciones,
el incremento es \(q(4,72)\), distinto de cero.
\end{proof}

\begin{remark}[Tipo de la extensión acumulativa]
La construcción anterior amplía historias dirigidas mediante contadores
positivos. No es un levantamiento automorfo del grupoide
\(\mathcal G_{\rm mem}\), ni demuestra que la coordenada ilimitada pertenezca
ya al estado \(\mathcal X_{\mathrm{TPK}}^{\mathrm{enr}}\). Su completación de Grothendieck
\(\mathbb Z^2\) puede utilizarse para análisis armónico, pero no transforma
el coste positivo
\(\mathcal A_\ast=G_{\rm or}+\lambda_\ast S\) en un carácter del grupo de
isotropía: al invertir formalmente una ruta, sus indicadores no adquieren
signo opuesto.
\end{remark}

\paragraph{Dominio de validez de la extensión acumulativa.}
La jerarquía \(27\to54\to108\), las incidencias de \(F_{\rm av}\) y los
cociclos anteriores son identidades exactas en la cronología y los
lectores declarados. La afirmación se restringe al estado observable que
esos lectores conservan; extenderla al estado enriquecido completo requeriría
probar el retorno de cada campo olvidado.

\section{Grupoide de memoria}
\label{sec:grupoide-memoria-tpk}

Sea \(\mathsf P_{\rm TPK}^{\rm enr}\) el grupoide de caminos enriquecidos:
sus objetos son estados TPK supervivientes y sus flechas son rutas admisibles
con origen y destino declarados, composición por concatenación e inversión
por reversión de ruta. Sea \(\sim_{\rm mem}\) una equivalencia que sólo
identifica dos flechas cuando conserva los campos de memoria requeridos por
el lector considerado y que es congruente con identidades, inversión y
concatenación. Definimos
\[
 \mathcal G_{\rm mem}
 :=\mathsf P_{\rm TPK}^{\rm enr}/\!\sim_{\rm mem}.
\]
Esta definición no transforma la memoria en torsión. La memoria pertenece al
estado discreto y distingue historias; la torsión pertenece a una conexión
geométrica y mide \(D_\omega e\). Sólo un mapa entre ambos dominios puede
relacionarlas.

Una realización geométrica exige, como datos adicionales, una longitud
\(\ell_0>0\), una representación del grupoide
\[
 \rho_{\rm C}:
 \mathcal G_{\rm mem}
 \longrightarrow
 \operatorname{Spin}(1,3)\ltimes\RR^{1,3},
\]
y un mapa de campos
\[
 \mathfrak R_{\rm grav}:
 \mathsf P_{\rm TPK}^{\rm enr}
 \dashrightarrow
 \mathsf{Cartan}_{1,3}.
\]
Para definir el segundo mapa deben especificarse su acción sobre objetos y
rutas, su compatibilidad con concatenación e inversión y la regularidad de la
conexión resultante. Ninguna de estas propiedades se deduce del mero conteo
de retornos.

\section{Conexión conjunta y condición de entrelazamiento}
\label{sec:conexion-cartan-conjunta}

Bajo las hipótesis geométricas del \cref{chap:gravedad} ---variedad
orientada de dimensión cuatro, co-tétrada \(e^I\), conexión de Lorentz
\(\omega^{IJ}\) y condiciones de borde allí declaradas---, la conexión de
Cartan conjunta se escribe
\[
 \mathcal A_{\ell_0}
 =
 \begin{pmatrix}
  \omega&\ell_0^{-1}e\\
  0&0
 \end{pmatrix}.
\]
Su curvatura es la identidad algebraica
\[
 \mathcal F_{\ell_0}
 =\dd\mathcal A_{\ell_0}
  +\mathcal A_{\ell_0}\wedge\mathcal A_{\ell_0}
 =
 \begin{pmatrix}
  F_\omega&\ell_0^{-1}T_{\rm tor}\\
  0&0
 \end{pmatrix},
\]
donde
\[
 F_\omega^{IJ}
 =\dd\omega^{IJ}+\omega^I{}_K\wedge\omega^{KJ},
 \qquad
 T_{\rm tor}^{I}=D_\omega e^I.
\]
Curvatura y torsión aparecen como componentes distintas de un mismo defecto
de transporte. La igualdad matricial se sigue algebraicamente de los datos
explícitos \(e,\omega,\ell_0\), pero no construye estas entradas desde TPK.

El contenido fuerte del puente es la conmutatividad de holonomías:
\[
 \boxed{
 \Hol_{\mathcal A_{\ell_0}}
 \bigl(\mathfrak R_{\rm grav}(\gamma)\bigr)
 =
 \rho_{\rm C}\!\left(
 \Hol_{\rm TPK}(\gamma)
 \right)}
 \qquad
 (\gamma\in\mathsf P_{\rm TPK}^{\rm enr}).
\]
La igualdad debe ser compatible con identidades, inversión, concatenación,
subdivisión, cambio de punto base y conjugación. Sólo entonces la memoria
discreta queda representada en una conexión; no queda rebautizada como
torsión ni como curvatura.

\begin{criterion}[Condiciones de existencia del puente TPK--Cartan]
\label{crit:puente-tpk-cartan}
No existe una realización con la conmutatividad y la compatibilidad declaradas
si ocurre cualquiera de los siguientes hechos:
\begin{enumerate}
 \item dos representantes equivalentes de una misma ruta producen
 holonomías geométricas no conjugadas;
 \item la imagen no preserva identidades, inversos o concatenación;
 \item una subdivisión admisible altera la holonomía;
 \item la igualdad de holonomías falla para \(\mathcal K_{27}\) o para sus
 concatenaciones de profundidades \(54\) y \(108\);
 \item el mapa identifica memoria y torsión sin exhibir un morfismo que
 preserve sus tipos;
 \item alguna elección descendente ---la acción de Holst, \(G\), una masa o
 una carta SI--- se utiliza para redefinir \(A\), \(C^\ast\) o el retorno
 TPK de origen.
\end{enumerate}
\end{criterion}

\paragraph{Antecedencia del funcional \(\Gamma_{6\to8}\) respecto de la realización Palatini--Cartan--Holst.}
El \cref{mdg:chap:barbero} fija, una vez construidos \(A,C^\ast,q_\pm\), el
funcional interno \(\Gamma_{6\to8}\). El \cref{chap:gravedad} estudia después
la acción Palatini--Cartan--Holst y sus ecuaciones variacionales. El puente
presente no introduce una ruta inversa entre ambas etapas.

\endgroup

% Owner íntegro de la realización discreta del retorno, la holonomía y la
% respuesta colectiva. Se sitúa después del puente que declara sus hipótesis.
\begin{HMTScientificOwnerSubordinated}
\chapter{Gravedad: holonomía, torsión y respuesta colectiva}
\label{chap:gravedad}
\index{gravedad}\index{Cartan--Holst}\index{holonomía}
\index{torsión}\index{gravitón!no primitivo}

La gravedad no se introduce aquí mediante una fórmula decimal para \(G\).
Su construcción reúne dos ramas ya tipadas. En la primera, una ruta conserva
orientación y memoria; una conexión compara marcos en puntos distintos; su
holonomía mide el retorno y su curvatura, el defecto de cierre. En la segunda,
la carta dimensional construye \(G_{\rm int}(\theta_G)\) y
\(\kappa_{\rm int}\) a partir de la escala interna de acción, la velocidad y
la duración elemental. Sólo después de disponer de ambas ramas se escribe la
acción Palatini--Cartan--Holst. Ninguna cifra metrológica externa participa en
esa precedencia.

\section{Retornos: posición, orientación, fase y memoria}

Sea \(F_{\rm obs}\) la cronología observable del transporte y
\[
 \mathcal K_{27}=F_{\rm obs}^{27}.
\]
La jerarquía exacta distingue tres retornos:
\[
 \begin{array}{c|c}
 27\ \text{pasos}&
 \text{retorno de posición e inversión de orientación},\\
 54\ \text{pasos}&
 \text{retorno de posición y orientación},\\
 108\ \text{pasos}&
 \text{retorno completo de la fase observable}.
 \end{array}
\]
El retorno observable no borra el registro acumulado.  Si \(g(e)\) cuenta
una inversión de orientación y \(s(e)\) un cambio efectivo de hoja, entonces
\[
 c(\gamma)=\sum_{e\subset\gamma}(g(e),s(e))
\]
satisface la ley cocíclica de concatenación.  En los bloques canónicos,
\[
 c_{27}=(1,18),\qquad c_{108}=(4,72).
\]
Por tanto, la proyección observable de \(\mathcal K_{27}\) tiene orden
cuatro, mientras la coordenada acumulativa sigue creciendo.  Volver a la
misma fase no significa volver a la misma historia.

\begin{theorem}[Cobertura acumulativa del retorno]
\label{thm:cobertura-retorno-grav-rev6}
Sobre \(\mathcal O_{108}\times\mathbb N^2\), la elevación
\[
 \widetilde F(x,m)=
 \bigl(F_{\rm obs}x,m+c(x\to F_{\rm obs}x)\bigr)
\]
satisface
\[
 \widetilde{\mathcal K}_{27}(x,m)
 =\bigl(\mathcal K_{27}x,m+(1,18)\bigr)
\]
y
\[
 \widetilde{\mathcal K}_{27}^{\,4}(x,m)
 =\bigl(x,m+(4,72)\bigr)\neq(x,m).
\]
\end{theorem}

\begin{proof}
La constancia de \(c_{27}\) sobre la órbita y la ley de concatenación dan
\[
 c_{108}=\sum_{j=0}^{3}c_{27}(\mathcal K_{27}^jx)
 =4(1,18).
\]
La primera coordenada retorna; la segunda conserva la memoria acumulada.
\end{proof}

Éste es el dato que una realización geométrica debe preservar.  No se
identifica la memoria discreta con la torsión: se construye una representación
que permita compararlas.

\section{Del grupoide de caminos al grupo de Cartan}

Sea \(\mathcal G_{\rm mem}\) el grupoide de rutas enriquecidas, con
composición por concatenación e inversa por reversión. Fijemos un elemento
\(R_4\in\operatorname{Spin}^+(1,3)\) de orden cuatro y un vector temporal
unitario \(u\in\mathbb R^{1,3}\) fijado por \(R_4\). Si
\[
c(\gamma)=\bigl(c_g(\gamma),c_s(\gamma)\bigr)\in\mathbb Z^2
\]
es el cociclo aditivo de inversión y cambio de hoja, definimos
\[
\rho_{\rm C}^{\rm disc}(\gamma)
=
\left(
R_4^{\,c_g(\gamma)},
\ell_0c_s(\gamma)u
\right)
\in\operatorname{Spin}^+(1,3)\ltimes\mathbb R^{1,3}.
\]

\begin{theorem}[Realización discreta del cociclo de retorno]
\label{thm:realizacion-discreta-cartan-rev6}
\(\rho_{\rm C}^{\rm disc}\) es un funtor del grupoide de rutas en el grupo de
Cartan. En particular, conserva identidad, concatenación e inversión, y
transporta el retorno \(27\to54\to108\) sin borrar la coordenada acumulativa.
\end{theorem}

\begin{proof}
La ley del producto semidirecto es
\[
(R,a)(S,b)=(RS,a+Rb).
\]
Como \(R_4u=u\) y \(c\) es aditivo,
\[
\begin{aligned}
\rho_{\rm C}^{\rm disc}(\gamma_2)
\rho_{\rm C}^{\rm disc}(\gamma_1)
&=
\left(
R_4^{c_g(\gamma_2)+c_g(\gamma_1)},
\ell_0\bigl(c_s(\gamma_2)+c_s(\gamma_1)\bigr)u
\right)\\
&=\rho_{\rm C}^{\rm disc}(\gamma_2\gamma_1).
\end{aligned}
\]
La ruta identidad tiene cociclo nulo y la reversión cambia su signo, por lo
que se preservan identidad e inversa. Las fórmulas de retorno siguen de
\(c_{27}=(1,18)\) y \(c_{108}=4c_{27}\).
\end{proof}

Este teorema construye el puente discreto. Una realización suave posterior
consiste en una representación
\[
 \rho_{\rm C}:\mathcal G_{\rm mem}
 \longrightarrow\operatorname{Spin}(1,3)\ltimes\mathbb R^{1,3}
\]
y en campos \(e^I,\omega^{IJ}\) cuya holonomía entrelace esa representación.
El diagrama que fija el tipo es
\[
 \operatorname{Hol}_{\mathcal A_{\ell_0}}
   \bigl(\mathfrak R_{\rm C}(\gamma)\bigr)
 =
 \rho_{\rm C}\bigl(\operatorname{Hol}_{\rm TPK}(\gamma)\bigr).
\]
Identidades, inversión, concatenación y subdivisión deben conservarse.  Esta
ecuación no rebautiza el acarreo como curvatura: exige que el transporte
discreto ya representado y el geométrico realicen la misma composición.

Fijada una escala positiva \(\ell_0\), la conexión conjunta es
\[
 \mathcal A_{\ell_0}
 =
 \begin{pmatrix}
 \omega&\ell_0^{-1}e\\
 0&0
 \end{pmatrix}.
\]

\begin{proposition}[Curvatura conjunta de Cartan]
\label{prop:curvatura-conjunta-cartan-rev6}
La curvatura de \(\mathcal A_{\ell_0}\) es
\[
 \mathcal F_{\ell_0}
 =\dd\mathcal A_{\ell_0}
  +\mathcal A_{\ell_0}\wedge\mathcal A_{\ell_0}
 =
 \begin{pmatrix}
 F_\omega&\ell_0^{-1}T_{\rm tor}\\
 0&0
 \end{pmatrix},
\]
donde
\[
 F_\omega^{IJ}
 =\dd\omega^{IJ}+\omega^I{}_K\wedge\omega^{KJ},
 \qquad
 T_{\rm tor}^{I}=D_\omega e^I.
\]
\end{proposition}

\begin{proof}
Se multiplican las matrices de formas usando el producto exterior.  El
bloque diagonal produce \(F_\omega\); el bloque de traslación produce
\(\dd e+\omega\wedge e=D_\omega e\).
\end{proof}

Curvatura y torsión son dos componentes de un mismo defecto de transporte,
pero cumplen funciones distintas.  La curvatura mide el retorno de marcos;
la torsión mide el cierre de la co-tétrada.

\section{Gauss--Stokes y capacidad areal}

En un complejo celular finito, para toda \(1\)-cocadena \(\Omega\) y toda
\(2\)-cadena \(z\),
\[
 \boxed{\quad
 \langle\delta\Omega,z\rangle
 =\langle\Omega,\partial z\rangle.
 \quad}
\]
La identidad es la definición del coborde como adjunto de la frontera.
Cuando \(z\) suma las caras coherentemente orientadas de una
pseudovariedad, las aristas interiores se cancelan y sólo permanece el
borde geométrico.  La misma igualdad se leyó ya como
vorticidad--circulación y como curvatura--holonomía.

La escala areal también posee un modelo exacto.  En el grafo cúbico
\(N\times N\times N\), todo corte entre dos caras opuestas contiene al menos
\(N^2\) aristas y una capa transversal alcanza la cota.  Para \(N=9\),
\[
 9^3=729=27^2.
\]
Esta igualdad cardinal relaciona volumen y corte mínimo; no identifica los
grupos \((\mathbb Z/9)^3\) y \((\mathbb Z/27)^2\), que tienen exponentes
distintos.

\section{La acción Palatini--Cartan--Holst}

Sea \(M\) una variedad lisa orientada de dimensión cuatro, con co-tétrada no
degenerada \(e^I\), conexión de Lorentz
\(\omega^{IJ}=-\omega^{JI}\), y condiciones de borde que anulen el término
variacional de frontera.  Escribamos
\[
 \Sigma^{IJ}=e^I\wedge e^J,\qquad
 (\star X)^{IJ}=\frac12\epsilon^{IJ}{}_{KL}X^{KL},
 \qquad
 G_{\rm int}=G_{\rm int}(\theta_G),\qquad
 \kappa_{\rm int}=\frac{8\pi G_{\rm int}}{c_{\rm int}^{4}}.
\]
Estas dos últimas cantidades no se definen mediante la acción: proceden de
la construcción constitutiva de la \cref{eq:G-interno}. El
\cref{thm:homogeneidad-pch} demuestra antes que la combinación siguiente
pertenece a la recta de acción.
La acción de primer orden es
\[
\boxed{
 S[e,\omega,\Psi]
 =\frac1{2\kappa_{\rm int}c_{\rm int}}\int_M
 \Sigma_{IJ}\wedge
 \left(\star F^{IJ}+\frac1\gamma F^{IJ}\right)
 -\frac{\Lambda_{\rm int}}{\kappa_{\rm int}c_{\rm int}}
  \int_M\operatorname{vol}_e
 +S_{\rm m}[e,\omega,\Psi].}
\]
El parámetro interno construido en el capítulo anterior especializa
\[
 \gamma=\Gamma_{6\to8}.
\]
La sustitución se efectúa después de obtener
\(\Gamma_{6\to8}\) desde \(A,C^\ast,q_\pm\) y la inclusión
hexada--octada; la acción no se usa para redefinir esas entradas.

\begin{theorem}[Ecuación de Cartan--Holst]
\label{thm:cartan-holst-rev6}
Si la corriente de espín se normaliza mediante
\[
 \delta_\omega S_{\rm m}
 =-\frac12\int_M\delta\omega^{IJ}\wedge\sigma_{IJ},
\]
la variación respecto de \(\omega\) produce
\[
 \boxed{\quad
 D_\omega
 \left[(\star+\gamma^{-1}I)\Sigma^{IJ}\right]
 =\kappa_{\rm int}c_{\rm int}\,\sigma^{IJ}.
 \quad}
\]
\end{theorem}

\begin{proof}
Se usa \(\delta F^{IJ}=D_\omega\delta\omega^{IJ}\), se integra por partes y
se iguala a cero el coeficiente de la variación arbitraria.  El término
cosmológico no depende de \(\omega\).
\end{proof}

En firma lorentziana, \(\star^2=-I\) sobre bivectores.  Para
\(\gamma\in\mathbb R\setminus\{0\}\),
\[
 (\star+\gamma^{-1}I)^{-1}
 =\frac{\gamma^2}{1+\gamma^2}
  (-\star+\gamma^{-1}I).
\]
Por ello, si la corriente de espín se anula,
\[
 T_{\rm tor}^I=0,\qquad
 \omega=\mathring\omega(e).
\]
Con materia espinorial, la torsión se determina algebraicamente por la
corriente.  No constituye, en la acción mínima, un grado propagante
independiente.

\section{Ecuación métrica, espín y balance geométrico}

Para fijar sin ambigüedad la normalización, definimos las fuentes medidas en
la recta de acción por
\[
 \delta_g S_{\rm m}
 =-\frac12\int_M\operatorname{vol}_e\,
 \mathcal T^{S}_{\mu\nu}\,\delta g^{\mu\nu}.
\]
Cuando las coordenadas de la tétrada tienen dimensión de longitud, escribimos
\[
 T^{\rm phys}_{\mu\nu}:=c_{\rm int}\mathcal T^{S}_{\mu\nu},
 \qquad
 U^{\rm phys,spin}_{\mu\nu}
 :=c_{\rm int}\mathcal U^{S,\rm spin}_{\mu\nu}.
\]
Tras eliminar algebraicamente la conexión, la ecuación métrica adopta las
dos formas equivalentes
\[
 \boxed{
 G_{\mu\nu}+\Lambda_{\rm int}g_{\mu\nu}
 =\kappa_{\rm int}c_{\rm int}
 \bigl(\mathcal T_{\mu\nu}^{S,\rm LC}
       +\mathcal U_{\mu\nu}^{S,\rm spin}\bigr)
 =\kappa_{\rm int}
 \bigl(T_{\mu\nu}^{\rm phys,LC}
       +U_{\mu\nu}^{\rm phys,spin}\bigr).}
\]
La contribución de espín es cuadrática en la corriente correspondiente y
depende del acoplamiento de materia; su coeficiente no se fija sin declarar
la acción fermiónica y las convenciones de dualidad.

Las identidades geométricas
\[
 D_\omega T_{\rm tor}^{I}=F^{I}{}_{J}\wedge e^{J},
 \qquad
 D_\omega F^{IJ}=0
\]
implican, después de eliminar la torsión,
\[
 \nabla^\mu\bigl(T_{\mu\nu}^{\rm phys,LC}
             +U_{\mu\nu}^{\rm phys,spin}\bigr)=0.
\]
Este balance de Bianchi--Noether es una condición de consistencia del
transductor gravitatorio: ningún término adicional puede incorporarse sin
proceder de una acción compatible o satisfacer el mismo balance.

\section{Correspondencia areal–volumétrica \(\pi/\varphi^2\)}

El calendario \((\Sigma,\Pi,\Pi)\) produce la matriz de incidencia
\[
 F_{\rm av}=\begin{pmatrix}0&1\\1&1\end{pmatrix}.
\]
Su medida de Parry da
\(\mu_{\rm ar}/\mu_{\rm vol}=\varphi^{-2}\), y el lector regional marca una
sola vez el cierre \(\pi_{\rm HMT}\).  El retorno marcado satisface
\[
 \Theta_n
 =\pi_{\rm HMT}\frac{F_{n-1}}{F_{n+1}}
 \longrightarrow
 \frac{\pi_{\rm HMT}}{\varphi^2}.
\]
Esta razón no es otra versión de \(G\). Es la correspondencia adimensional entre
lectura areal y lectura volumétrica.  Bajo una misma amplitud positiva
\(\mathcal M\),
\[
 \mathcal I_{\rm av}=\mu_{\rm vol}\mathcal M,\qquad
 \mathcal G_{\rm av}=\pi_{\rm HMT}\mu_{\rm ar}\mathcal M
\]
dan exactamente
\[
 \frac{\mathcal G_{\rm av}}{\mathcal I_{\rm av}}
 =\frac{\pi_{\rm HMT}}{\varphi^2}.
\]
En la hoja plana, la restricción diagonal es
\(\mathcal G_{\rm tor}=\mathcal I_{\rm tor}\).  La equivalencia local y la
ambivalencia de frontera son, por tanto, dos regímenes del lector, no dos
valores incompatibles de una masa.

\section{Tipado de \(G\) como transductor gravitatorio y orden causal de sus lectores}

El tipo y la homogeneidad de \(G\) pertenecen a la construcción
constitutiva previa. La \cref{eq:G-interno} fija una sola vez la familia
covariante \(G_{\rm int}(\theta_G)\), y el
\cref{thm:homogeneidad-pch} demuestra que su inserción en la acción posee
dimensión de acción. Este capítulo no redefine esa familia: utiliza su
sección junto con la holonomía ya construida. La causalidad no es lineal,
sino convergente:
\[
\begin{aligned}
 \mathcal B_{\rm geom}&:
 \text{retorno y cociclo}\to\text{representación de Cartan}
 \to(e,\omega),\\
 \mathcal B_{\rm dim}&:
 (\hbar_{\rm int},c_{\rm int},t_0,\theta_G)
 \to G_{\rm int}\to\kappa_{\rm int},
\end{aligned}
\]
\[
 (\mathcal B_{\rm geom},\mathcal B_{\rm dim})
 \to S_{\rm PCH}
 \to\text{torsión algebraica, ecuación métrica y balance}.
\]
La comparación metrológica pertenece al final y no selecciona ninguna de
las dos ramas.

\Needspace{22\baselineskip}
\section{Origen colectivo del modo gravitatorio en la holonomía de rutas}

La acción anterior tiene como objetos básicos la co-tétrada y la conexión.
La torsión de Cartan es algebraica; la curvatura posee el sector propagante que
contiene las ondas gravitatorias.  Por ello el formalismo no necesita
postular un gravitón elemental para definir la interacción.  Una eventual
cuantización de las perturbaciones puede producir un modo colectivo de
helicidades \(\pm2\); ese modo sería una excitación del transporte
geométrico, no un dato añadido al nivel de APP, TRIT o TPK.

Esta distinción también excluye una confusión frecuente.  La existencia de
ondas clásicas demuestra propagación de curvatura; no demuestra por sí sola
un cuanto fundamental.  Recíprocamente, describir el campo mediante
geometría no elimina sus grados radiativos.  HMT--MD sitúa la pregunta en el
lugar correcto: primero se construye la memoria y su holonomía; después se
estudia el espectro de sus excitaciones.

\end{HMTScientificOwnerSubordinated}

\paragraph{Comparación tipada del cuanto de área con energía, entropía e información de horizonte.} La acción de primer orden permite comparar el cuanto areal con energía, entropía e información de horizonte sin fundir esos observables.
